\section{Unicritical multiplier curves and the reduced multiplier}
\label{sec:unicritical-reduction}

Fix integers $d\ge2$ and $m\ge1$, and put
\[
 f_c(z)=z^d+c,\qquad e=d-1,\qquad
 \nu_d(m)=\sum_{r\mid m}\mu(m/r)d^r.
\]
The $m$-th multiplier polynomial $M_{d,m}(c,x)$ is the monic polynomial
whose roots, for generic $c$, are the multipliers of the cycles of
exact period $m$.  It has degree $\nu_d(m)/m$ in $x$.  We use the
following standard facts about this polynomial and the corresponding
dynatomic curve:
\begin{enumerate}[label=\textup{(\alph*)}]
\item the marked dynatomic curve is geometrically irreducible;
\item $M_{d,m}$ is geometrically irreducible;
\item there is a polynomial $D_{d,m}(t,x)\in\Z[t,x]$ such that
      \[
       M_{d,m}(c,x)=D_{d,m}(d^dc^e,x),
       \numberthis\label{eq:scaled-multiplier}
      \]
      and $D_{d,m}$ has degree $\nu_d(m)/d$ in $t$, with leading
      coefficient $\pm1$.
\end{enumerate}
The geometric irreducibility statements go back to Morton
\cite[Corollary 1]{Morton1996}; the integrality, degree, and leading
coefficient in the scaled parameter are proved in
\cite[Theorem 1.2 and Propositions 3.3, 3.5]
{MurakamiSanoTakehira2024}.  See also
\cite[Theorem 1.38 and Propositions 1.51, 1.59]{Huguin2021Thesis}.

Let $\mathcal Y_{d,m}$ be the normalization of the marked dynatomic
curve, and write
\[
 L=\overline\Q(\mathcal Y_{d,m}).
\]
Iteration induces an automorphism $\sigma$ of order $m$.  If $z$ is
the marked point, put
\[
 z_j=f_c^j(z),\qquad
 P=\prod_{j=0}^{m-1}z_j.
 \numberthis\label{eq:cycle-product}
\]
Then $P\in L^{\langle\sigma\rangle}$ and the cycle multiplier is
\[
 x=(f_c^m)'(z)=d^mP^e.
 \numberthis\label{eq:multiplier-product}
\]

\begin{lemma}[The cycle and scaling quotients]
\label{lem:cycle-scaling-quotients}
One has
\[
 L^{\langle\sigma\rangle}=\overline\Q(c,x).
 \numberthis\label{eq:cycle-fixed-field}
\]
The action
\[
 \xi\cdot(c,z_0,\ldots,z_{m-1})
   =(\xi c,\xi z_0,\ldots,\xi z_{m-1}),
 \qquad \xi\in\mu_e,
 \numberthis\label{eq:scaling-action}
\]
commutes with $\sigma$, fixes $x$, and has fixed field
\[
 \overline\Q(c,x)^{\mu_e}
 =\overline\Q(c^e,x)=\overline\Q(t,x),
 \qquad t=d^dc^e.
 \numberthis\label{eq:scaling-fixed-field}
\]
\end{lemma}

\begin{proof}
The inclusion $\overline\Q(c,x)\subseteq L^{\langle\sigma\rangle}$
follows from \eqref{eq:multiplier-product}.  The two extensions over
$\overline\Q(c)$ have the same degree:
\[
 [L^{\langle\sigma\rangle}:\overline\Q(c)]
 =\frac{\nu_d(m)}m
 =\deg_x M_{d,m}.
\]
Here geometric irreducibility of $M_{d,m}$ is essential: it identifies
$\overline\Q(c,x)$ with the function field of the multiplier curve and
prevents a hidden proper subfield.  This proves
\eqref{eq:cycle-fixed-field}.

Since $\xi^d=\xi$, one has
$f_{\xi c}(\xi z)=\xi f_c(z)$, proving that
\eqref{eq:scaling-action} is an action.  It sends
$P$ to $\xi^mP$ and hence fixes $x$.  Its action on the transcendental
element $c$ is faithful, so the fixed field is
$\overline\Q(c^e,x)$, which is also $\overline\Q(t,x)$.
\end{proof}

Let $\mathcal X_{d,m}$ denote the normalization of the projective
closure of $D_{d,m}(t,x)=0$.  Thus
\[
 \Q(\mathcal X_{d,m})=\Q(t,x).
\]
Set
\[
 \gamma=(m,e),\qquad q=e/\gamma.
 \numberthis\label{eq:gq-definition}
\]

\begin{proposition}[The reduced multiplier]
\label{prop:reduced-multiplier}
The function
\[
 y=d^{m/\gamma}P^q
 \numberthis\label{eq:reduced-multiplier}
\]
belongs to $\Q(t,x)$ and satisfies
\[
 x=y^\gamma.
 \numberthis\label{eq:x-yg}
\]
Moreover,
\[
 \deg(y:\mathcal X_{d,m}\longrightarrow\mathbf P^1)
 =N_{d,m}:=\frac{\nu_d(m)}{d\gamma}.
 \numberthis\label{eq:reduced-degree}
\]
\end{proposition}

\begin{proof}
Because $mq$ is divisible by $e$, the expression in
\eqref{eq:reduced-multiplier} is fixed by $\mu_e$.  It is defined over
$\Q$, so Galois descent and Lemma~\ref{lem:cycle-scaling-quotients}
give $y\in\Q(t,x)$.  Equation \eqref{eq:x-yg} follows at once from
\eqref{eq:multiplier-product}.  Finally, the degree of the map $x$ on
$\mathcal X_{d,m}$ is $\deg_tD_{d,m}=\nu_d(m)/d$.  Since $x=y^\gamma$,
degrees of finite maps give \eqref{eq:reduced-degree}.
\end{proof}

The next theorem is the geometric core of the generalization.  The
stabilizer $h$ appearing in its proof cannot in general be omitted.

\begin{theorem}[Maximal isogeny and pullback irreducibility]
\label{thm:unicritical-pb}
The reduced multiplier
\[
 y:\mathcal X_{d,m}\dashrightarrow\mathbf G_m
\]
satisfies \emph{(PB)}.  Furthermore, $\gamma=(m,d-1)$ is the largest
geometric isogeny degree of the multiplier map $x$: if
\[
 x=w^r\qquad\text{in }\overline\Q(\mathcal X_{d,m})
\]
for some $r\ge1$, then $r\mid\gamma$.
\end{theorem}

\begin{proof}
We first suppose $m>1$.  Write $\Phi^*_{d,m}(z,c)$ for the $m$-th
dynatomic polynomial of $f_c$.  We first verify directly that its
critical specialization is squarefree.  Put
\[
 a_n(c)=f_c^{\circ n}(0)\qquad(n\geq1).
\]
Then $a_1=c$ and, for $n\geq2$, one has $a_n=a_{n-1}^d+c$, so
$a_n\in\Z[c]$ is monic and
\[
 a_n'=d a_{n-1}^{d-1}a_{n-1}'+1.
\]
If $\beta$ were a common root of $a_n$ and $a_n'$, then $\beta$ would be
an algebraic integer, as would
$a_{n-1}(\beta)^{d-1}a_{n-1}'(\beta)$.  The displayed identity would give
\[
 a_{n-1}(\beta)^{d-1}a_{n-1}'(\beta)=-\frac1d,
\]
which is impossible because $d\geq2$ and the rational algebraic
integers are the integers.  Thus every $a_n$ is squarefree.  The
dynatomic product identity, specialized at $z=0$, is
\[
 a_n(c)=\prod_{r\mid n}\Phi^*_{d,r}(0,c).
\]
Consequently every factor $\Phi^*_{d,r}(0,c)$ is squarefree.  Moreover,
M\"obius inversion applied to degrees gives
\[
 \deg_c\Phi^*_{d,m}(0,c)
 =\sum_{r\mid m}\mu(m/r)d^{r-1}
 =\frac{\nu_d(m)}d>0.
\]
Compare \cite[Proposition 1.58]{Huguin2021Thesis}.  Choose a root $c_0$
of $\Phi^*_{d,m}(0,c)$.  Because the
multiplier of the critical orbit is zero, the formal-period criterion
shows that $0$ has exact period $m$ for $f_{c_0}$; in particular,
$c_0\ne0$.  Squarefreeness gives
\[
 \partial_c\Phi^*_{d,m}(0,c_0)\ne0.
\]
Thus the marked dynatomic curve is smooth at $(0,c_0)$ and $z=z_0$ is
a local parameter there.  The remaining points $z_1,\ldots,z_{m-1}$
of the critical cycle are nonzero, so $P=z_0\cdots z_{m-1}$ has a
simple zero.  Exact periodicity makes the cycle-rotation stabilizer
trivial, and the scaling action is free because $c_0\ne0$.
Consequently both quotient maps are unramified at this point.  At its
image $P_0$ on $\mathcal X_{d,m}$ one therefore has
\[
 \ord_{P_0}(y)=q,\qquad \ord_{P_0}(x)=e.
 \numberthis\label{eq:center-orders}
\]

We next compute an order at infinity.  Let
$\overline{\mathcal Y}_{d,m}$ be the smooth projective model of the
marked dynatomic curve, choose a point $R_\infty$ above $c=\infty$, and
normalize the base change given by $c=s^{-d}$.  Choose a point
$\widetilde R_\infty$ above $(R_\infty,s=0)$, write
$v=\ord_{\widetilde R_\infty}$, and put
\[
 b_j=\frac{v(z_j)}{v(s)}.
\]
The recurrence $z_{j+1}=z_j^d+c$ forces $b_j=-1$ for every $j$.
Indeed, if the minimum of the $b_j$ were less than $-1$, then at an
index where it is attained the term $z_j^d$ would have strictly smaller
valuation than $c$, giving $b_{j+1}=db_j<b_j$, a contradiction.  Thus
all $b_j\ge-1$.  If some $b_j>-1$, then $c$ would be the unique term of
least valuation and would give $b_{j+1}=-d<-1$, another contradiction.

It follows that $z_j=s^{-1}u_j$ with every $u_j$ a unit.  The orbit
equations become
\[
 u_j^d+1=s^e u_{j+1}\qquad(j\bmod m).
 \numberthis\label{eq:infinity-equations}
\]
Write $a_j=u_j(0)$, so $a_j^d=-1$.  At $s=0$ the Jacobian with respect
to $u_0,\ldots,u_{m-1}$ is the invertible diagonal matrix
\[
 \operatorname{diag}(da_0^{d-1},\ldots,da_{m-1}^{d-1}).
\]
The formal implicit-function theorem therefore identifies the completed
local ring at $\widetilde R_\infty$ with $\overline\Q[[s]]$.  Uniqueness
of the solution, and the fact that the equations depend on $s$ only
through $s^e$, also give
\[
 u_j(s)\in\overline\Q[[s^e]]^\times.
\]

The deck group $\mu_d$ of the base change acts by
\[
 (s,u_0,\ldots,u_{m-1})
 \longmapsto(\omega s,\omega u_0,\ldots,\omega u_{m-1}),
 \qquad \omega\in\mu_d.
\]
It leaves the original functions $c$ and $z_j=s^{-1}u_j$ unchanged.
At $s=0$ it sends $(a_j)_j$ to $(\omega a_j)_j$, so no nonidentity
element fixes $\widetilde R_\infty$.  The quotient of the normalized
base change by this deck action is $\overline{\mathcal Y}_{d,m}$;
hence the map back to the marked curve is unramified at
$\widetilde R_\infty$.  From
\[
 c=s^{-d},\qquad
 P=s^{-m}\prod_{j=0}^{m-1}u_j(s)
\]
we obtain on the marked curve
\[
 \ord_{R_\infty}(c)=-d,\qquad
 \ord_{R_\infty}(P)=-m.
 \numberthis\label{eq:marked-infinity-orders}
\]

Let $h$ be the order of the stabilizer of $R_\infty$ under cycle
rotation.  In characteristic zero the marked-to-cycle quotient has
ramification index $h$ at this point.  Since both $c$ and $P$ descend,
at the image $Q_\infty$ one has
\[
 \ord_{Q_\infty}(c)=-\frac dh,\qquad
 \ord_{Q_\infty}(P)=-\frac mh.
 \numberthis\label{eq:cycle-infinity-order}
\]
In particular, $h\mid d$ and $h\mid m$.

Every element of $\mu_e$ fixes $Q_\infty$.  On the normalized base
change the scaling action by $\xi\in\mu_e$ lifts locally as
\[
 (s,u_0,\ldots,u_{m-1})
 \longmapsto(\xi^{-1}s,u_0,\ldots,u_{m-1}).
\]
This covers $(c,z_j)\mapsto(\xi c,\xi z_j)$ and preserves the chosen
branch because $u_j(s)\in\overline\Q[[s^e]]$.  Thus the
induced action on the cycle quotient is faithful, since $\xi$ sends
the nonconstant function $c$ to $\xi c$.  Hence the stabilizer of
$Q_\infty$ is all of $\mu_e$, and in characteristic zero the
cycle-to-scaling quotient has ramification index exactly $e$ at
$Q_\infty$.
Since $y=d^{m/\gamma}P^q$ and $x=d^mP^e$, their orders at the image
$P_\infty$ on $\mathcal X_{d,m}$ are
\[
 \ord_{P_\infty}(y)=-\frac{m}{\gamma h},
 \qquad
 \ord_{P_\infty}(x)=-\frac mh.
 \numberthis\label{eq:infinity-orders}
\]

Write $m=\gamma r$.  Since $e=\gamma q$ and $(q,r)=1$, while $h\mid d$
and $\gamma\mid d-1$, one has $(h,\gamma)=1$.  The divisibility
$h\mid m=\gamma r$ therefore gives $h\mid r$, and hence
\[
 \gcd\left(q,\frac{m}{\gamma h}\right)
 =\gcd\left(q,\frac rh\right)=1.
 \numberthis\label{eq:coprime-orders}
\]
The two orders in \eqref{eq:center-orders} and
\eqref{eq:infinity-orders} show that $y$ is not a $p$-th power in
$\overline\Q(\mathcal X_{d,m})$ for any prime $p$.  The binomial
criterion \cite[Chapter VI, \S9, Theorem 9.1]{Lang2002} now makes
$U^n-y$ irreducible for every $n\ge1$.  When $4\mid n$, the exceptional
condition $y\in-4\overline\Q(\mathcal X_{d,m})^4$ is already excluded
because $-4=(1+i)^4$ in the constant field.  This proves (PB).

The two displayed orders of $x$ have greatest common divisor
\[
 \gcd\left(e,\frac mh\right)
 =\gamma\gcd\left(q,\frac rh\right)=\gamma.
\]
If $x=w^\ell$, every divisor order of $x$ is divisible by $\ell$, so
$\ell\mid\gamma$.  Since $x=y^\gamma$, the bound is attained.

It remains to treat $m=1$.  If $z$ is a fixed point, then
$c=z-z^d$ and $x=dz^{d-1}$.  Direct elimination gives
\[
 t=x(d-x)^{d-1}.
 \numberthis\label{eq:period-one-curve}
\]
Here $\gamma=1$ and $y=x$; in particular $y$ has order one at the center.
It satisfies (PB), and the maximal isogeny degree is $1$.  Notice that
the center is fully ramified in the scaling quotient in this case,
which is why it was separated from the preceding argument.
\end{proof}

\section{Torsion separation on the unicritical quotient}
\label{sec:unicritical-separation}

To apply Theorem~\ref{thm:kummer-torsion}, it remains to recover the
reduced multiplier from the scaled parameter at torsion fibers.

\begin{proposition}[Torsion separation]
\label{prop:unicritical-separation}
On the union of the finite fibers
\[
 y^{-1}(\theta),\qquad
 \theta\in\mu_\infty,\quad \theta^\gamma\ne1,
\]
the function $t=d^dc^{d-1}$ is injective.  In particular, if
$Q\in y^{-1}(\theta)$ and $t_0=t(Q)$, then
\[
 \theta\in\Q(t_0).
\]
\end{proposition}

\begin{proof}
First, a formal-period-$m$ point with $x\ne1$ has exact period $m$.
Indeed, if its exact period were $r<m$, the formal-period criterion
\cite[Theorem 4.5(b)]{Silverman2007} would write $m=rs$, with the
$r$-cycle multiplier a primitive $s$-th root of unity; its multiplier
for the $m$-th iterate would then be $1$.

A unicritical polynomial has at most one finite cycle whose multiplier
is a root of unity.  Indeed, after passing to a suitable iterate, such
a cycle is parabolic with multiplier $1$.  The Fatou critical-point
theorem for parabolic basins \cite[Theorem 10.15]{Milnor2006} supplies
a critical point of that iterate whose orbit tends to the cycle.  By
the chain rule, one of its forward images is a critical point of the
original polynomial, so the cycle attracts a critical orbit of
$f_c$.  The resulting critical point of $f_c$ is finite, since infinity
lies in the basin of the superattracting fixed point at infinity.  The
only finite critical point of $z^d+c$ is $0$, and the attracting basins
of two distinct periodic cycles are disjoint, so two such cycles cannot
coexist.  Also
$c=0$ cannot occur here: for $z^d$, a nonzero periodic orbit has
multiplier a positive power of $d$, while the zero fixed point has
multiplier zero.

Suppose two points $Q_1,Q_2$ in the indicated torsion fibers have the
same $t$-coordinate, represented by parameters $c_1,c_2$.  Since these
parameters are nonzero, there is a unique $\xi\in\mu_e$ with
$c_2=\xi c_1$.  The conjugacy $z\mapsto\xi z$ transports the first
cycle to a cycle for $f_{c_2}$.  The transported cycle and the second
cycle coincide by the preceding paragraph.  Moreover $P\mapsto\xi^mP$,
whereas
\[
 \xi^{mq}=1,
\]
so the value of $y=d^{m/\gamma}P^q$ is unchanged.  Thus $Q_1=Q_2$.

There is no hidden splitting on normalization: at these points
$\partial_z(f_c^m(z)-z)=x-1\ne0$, cycle rotation is free, and the
$\mu_e$-action is free because $c\ne0$.  Both quotient maps are \'etale
there.  The final assertion now follows either directly or from the
Galois argument used in Theorem~\ref{thm:kummer-torsion}.
\end{proof}

\section{Eventual exact factorization in the scaled parameter}
\label{sec:unicritical-factorization}

Let $A_{d,m}(T,Y)$ be the monic minimal polynomial of $t$ over
$\Q(y)$.  It has degree $N_{d,m}$ in $T$.  Since
$D_{d,m}(T,Y^\gamma)$ is monic up to a fixed sign in $T$, Gauss's lemma and
normality of $\Q[Y]$ allow $A_{d,m}$ to be normalized in
$\Q[Y,T]$.  Over $\overline\Q$ one has
\[
 D_{d,m}(T,Y^\gamma)
 \doteq
 \prod_{\xi\in\mu_\gamma}A_{d,m}(T,\xi Y),
 \numberthis\label{eq:generic-kummer-factorization}
\]
where $\doteq$ denotes equality up to the fixed leading sign.  Each
$A_{d,m}(T,\xi Y)$ is irreducible and divides the left-hand side.  The
factors are distinct: otherwise
$A_{d,m}(T,Y)=A_{d,m}(T,\delta Y)$ for some nontrivial
$\delta\in\mu_\gamma$.  Then $t\mapsto t$ and $y\mapsto\delta y$ would
induce a nontrivial automorphism of
$\overline\Q(t,y)=\overline\Q(\mathcal X_{d,m})$ fixing both $t$ and
$x=y^\gamma$.  This contradicts
$\overline\Q(\mathcal X_{d,m})=\overline\Q(t,x)$.  The factors are
therefore pairwise coprime, and comparison of degrees in $T$ gives
\eqref{eq:generic-kummer-factorization}.

For $k\ge2$, define the scaled unicritical delta polynomial by taking
the positive-leading normalization of
\[
 \mathcal P_{d,m,k}(T)
 \doteq\Res_X\bigl(\Phi_k(X),D_{d,m}(T,X)\bigr).
 \numberthis\label{eq:scaled-delta}
\]
In the standard unicritical notation of
\cite[Theorem 1.1]{MurakamiSanoTakehira2024},
\[
 \mathcal P_{d,m,k}(d^dc^{d-1})\doteq\Delta_{mk,m}(c),
\]
where $\doteq$ accounts only for the chosen leading-sign
normalization.  By
\eqref{eq:generic-kummer-factorization},
\[
 \mathcal P_{d,m,k}(T)
 =\prod_{\ord(\theta^\gamma)=k}A_{d,m}(T,\theta).
 \numberthis\label{eq:scaled-delta-reduced-product}
\]

\begin{theorem}[Eventual exact factorization]
\label{thm:unicritical-eventual-factorization}
Fix $d\ge2$ and $m\ge1$, and retain
\[
 \gamma=(m,d-1),\qquad
 N_{d,m}=\frac{\nu_d(m)}{d\gamma}.
\]
For all sufficiently large $k\ge2$:
\begin{enumerate}[label=\textup{(\roman*)}]
\item over $\Q^{\mathrm{ab}}$, $\mathcal P_{d,m,k}$ has exactly
      $\gamma\varphi(k)$ distinct irreducible factors, each of degree
      $N_{d,m}$;
\item over $\Q$, its complete irreducible factorization is indexed by
      \[
       a\mid\gamma,\qquad (k,\gamma/a)=1,
      \]
      and the factor indexed by $a$ has degree
      $N_{d,m}\varphi(ka)$;
\item if $t_0$ belongs to the factor indexed by a reduced multiplier
      $\theta$ of order $ka$, then
      \[
       \Q(t_0)\cap\Q^{\mathrm{ab}}=\Q(\theta).
       \numberthis\label{eq:unicritical-field-intersection}
      \]
\end{enumerate}
The corresponding factorization over any fixed number field $B$ is
given by Theorem~\ref{thm:kummer-torsion-number-field}.
\end{theorem}

\begin{proof}
Theorem~\ref{thm:unicritical-pb} supplies (PB), and
Proposition~\ref{prop:unicritical-separation} supplies torsion
separation for every $x\ne1$.  Since $k\ge2$, all fibers in
\eqref{eq:scaled-delta-reduced-product} have $x\ne1$.  Apply
Theorem~\ref{thm:kummer-torsion}.
\end{proof}

\begin{corollary}
\label{cor:unicritical-eventual-criterion}
For all sufficiently large $k$,
\[
 \mathcal P_{d,m,k}\text{ is irreducible over }\Q
 \quad\Longleftrightarrow\quad
 \operatorname{rad}((m,d-1))\mid k.
 \numberthis\label{eq:unicritical-radical-criterion}
\]
If $(m,d-1)=1$, the whole row is therefore eventually irreducible.
For the quadratic family $d=2$, this holds for every fixed $m$.
Thus every fixed row of the Morton--Vivaldi conjecture has only
finitely many possible exceptions.
\end{corollary}

\begin{proof}
This is Corollary~\ref{cor:abstract-factor-count}.  When $d=2$, the
scaled parameter $t=4c$ is an invertible linear change of the original
parameter, so no extra composition issue occurs.
\end{proof}

\begin{remark}[Effectivity]
Zannier notes that the arguments used for the toric theorem are
effective \cite[p.~400]{Zannier2010}.  Hence, for each fixed pair
$(d,m)$ and fixed base number field, the finite exceptional set is
effectively computable in principle.  No practical numerical bound is
claimed here.
\end{remark}
